\chapter{配置系统}

CCB 配置文件是 \src{.ccb/ccb.config}。源码支持三种文档形态：compact、rich TOML 和 hybrid。所有形态最终都被转换和验证为 version 2 project config。

\section{全局规则}

当前 CCB 项目配置的权威目标是 \cfg{version = 2}。rich TOML 和 windows topology 示例应显式写出 \cfg{version = 2}；compact 配置虽然不直接写这个字段，但 loader 会把它规范化为 version 2 project config。

\section{配置来源优先级}

CCB 按以下顺序加载配置：

\begin{enumerate}
  \item 项目配置：\src{.ccb/ccb.config}。
  \item 用户默认配置。
  \item 内建默认配置。
\end{enumerate}

如果项目中存在 \src{.ccb/ccb.config}，它优先。

\section{compact 配置}

compact 配置适合快速声明 agent 布局。每个 agent leaf 需要写成 \cfg{agent_name:provider}。

\begin{lstlisting}
cmd | archi:codex | reviewer:claude
\end{lstlisting}

规则：

\begin{itemize}
  \item \cfg{cmd} 是保留 token，表示前台命令 pane。
  \item \cfg{cmd} 不能写 provider。
  \item \cfg{cmd} 不能重复。
  \item 普通 agent leaf 必须声明 provider。
  \item compact config 至少要定义一个 agent。
  \item compact 中可以使用 role id shorthand，例如 \cfg{agentroles.archi:codex}。
\end{itemize}

compact 配置会生成：

\begin{itemize}
  \item \cfg{version = 2}
  \item \cfg{default_agents}
  \item \cfg{agents}
  \item \cfg{cmd_enabled}
  \item \cfg{layout}
\end{itemize}

\section{rich TOML}

rich TOML 适合完整配置。基本形态：

\begin{lstlisting}
version = 2
default_agents = ["agent1"]
cmd_enabled = true
layout = "cmd | agent1"

[agents.agent1]
provider = "codex"
target = "."
workspace_mode = "inplace"
restore = "auto"
permission = "manual"
\end{lstlisting}

rich TOML 顶层允许：

\begin{itemize}
  \item \cfg{version}
  \item \cfg{default_agents}
  \item \cfg{agents}
  \item \cfg{cmd_enabled}
  \item \cfg{layout}
  \item \cfg{ui}
  \item \cfg{windows}
  \item \cfg{tool_windows}
  \item \cfg{entry_window}
  \item \cfg{maintenance}
\end{itemize}

\section{hybrid 配置}

hybrid 是 compact layout 后面接 TOML overlay：

\begin{lstlisting}
cmd | archi:codex | reviewer:claude

[agents.archi]
model = "gpt-5"
labels = ["architecture"]

[maintenance.heartbeat]
enabled = true
\end{lstlisting}

hybrid overlay 只允许 \cfg{agents} 和 \cfg{maintenance}。它不能定义 compact layout 之外的新 agent，也不能覆盖 compact header 已经拥有的 \cfg{provider} 或 \cfg{workspace_mode}。

\section{基础布局模式}

没有 \cfg{windows} 时是基础布局模式。这是最简单、最常见的配置入口，适合一个前台 cmd pane 加若干 agent pane：

\begin{lstlisting}
version = 2
default_agents = ["archi", "reviewer"]
cmd_enabled = true
layout = "cmd | archi | reviewer"

[agents.archi]
provider = "codex"

[agents.reviewer]
provider = "claude"
\end{lstlisting}

规则：

\begin{itemize}
  \item \cfg{default_agents} 必需。
  \item 可以使用 \cfg{layout}。
  \item 可以使用 \cfg{cmd_enabled}。
  \item 不能使用 \cfg{ui.sidebar}。
  \item 不能使用 \cfg{entry_window}。
  \item \cfg{tool_windows} 需要 windows topology，因此基础布局模式不能用。
\end{itemize}

\section{windows topology 模式}

声明 \cfg{windows} 后进入 windows topology 模式：

\begin{lstlisting}
version = 2
entry_window = "main"

[windows]
main = "archi:codex | reviewer:claude"
ops = "ops:codex"
\end{lstlisting}

规则：

\begin{itemize}
  \item \cfg{windows} 不能为空。
  \item 每个 window leaf 必须声明 provider。
  \item 不支持 \cfg{cmd} leaf。
  \item agent 不能跨 windows 重复。
  \item \cfg{default_agents} 自动从 window leaves 推导，不能手写。
  \item \cfg{layout} 和 \cfg{cmd_enabled} 禁止使用。
  \item \cfg{entry_window}、\cfg{ui.sidebar}、\cfg{tool_windows} 可用。
\end{itemize}

\section{tool windows}

tool window 是非 agent 工具窗口：

\begin{lstlisting}
[tool_windows.files]
command = "CCB_WORKBENCH_PROFILE=rich CCB_WORKBENCH_FORCE_RICH=1 ccb-workbench files"
label = "files"
show_in_sidebar = true
\end{lstlisting}

规则：

\begin{itemize}
  \item \cfg{tool_windows} 需要 windows topology。
  \item 每个 tool window 必须有 \cfg{command}。
  \item 内置文件/workbench 能力首选在 \cfg{windows} 中写 \src{rich} alias，例如 \src{main = "agent1:codex, rich"}。
  \item \cfg{label} 可选。
  \item \cfg{show_in_sidebar} 可选，默认 true。
\end{itemize}

\section{sidebar}

sidebar 配置示例：

\begin{lstlisting}
[ui.sidebar]
mode = "every_window"
width = "15%"
bottom_height = 20
position = "right"
agents_height = "50%"
comms_height = "15%"
tips_height = "35%"
comms_limit = 5
comms_compact = true
tips_enabled = true
tips = ["Run ccb doctor when status is degraded."]
\end{lstlisting}

\cfg{position} 可选 \cfg{left} 或 \cfg{right}，默认 \cfg{left}。当前 sidebar 仍是纵向 pane；底部横排布局不属于该配置。旧的 \cfg{ui.sidebar.view} 子表仍可读取，但新配置建议把 sidebar 字段放在同一个 \cfg{ui.sidebar} 表中。

\section{agent spec}

agent 字段包括：

\begin{longtable}{p{0.30\linewidth}p{0.60\linewidth}}
\toprule
字段 & 用途 \\
\midrule
\cfg{provider} & 必需，provider 名。 \\
\cfg{target} & 工作目标，常见为 \cfg{"."}。 \\
\cfg{workspace_mode} & \cfg{inplace} 或 \cfg{git-worktree}。 \\
\cfg{workspace_root} & 外部 workspace root。 \\
\cfg{workspace_path} & 固定 workspace path。 \\
\cfg{workspace_group} & 共享 worktree group。 \\
\cfg{provider_command_template} & provider 启动命令模板。 \\
\cfg{runtime_mode} & runtime 类型，默认 pane-backed。 \\
\cfg{restore} & 恢复策略。 \\
\cfg{permission} & 权限策略。 \\
\cfg{queue_policy} & 队列策略，默认 serial-per-agent。 \\
\cfg{model} & provider model。 \\
\cfg{key}、\cfg{url} & API shortcut。 \\
\cfg{startup_args} & provider 启动参数。 \\
\cfg{env} & agent 环境变量。 \\
\cfg{api} & API 配置。 \\
\cfg{provider_profile} & provider profile。 \\
\cfg{branch_template} & worktree branch 模板。 \\
\cfg{labels} & 标签。 \\
\cfg{description} & 描述。 \\
\cfg{role} & role id。 \\
\cfg{watch_paths} & 关注路径。 \\
\bottomrule
\end{longtable}

\section{provider profile 和 API shortcut}

provider profile 字段包括 \cfg{mode}、\cfg{home}、\cfg{env}、\cfg{inherit_api}、\cfg{inherit_auth}、\cfg{inherit_config}、\cfg{inherit_skills}、\cfg{inherit_commands}、\cfg{inherit_memory}。

\cfg{key} 和 \cfg{url} shortcut 会转成 provider API env。注意：

\begin{itemize}
  \item shortcut 只支持特定 provider。
  \item 不能和 \cfg{env} 或 \cfg{provider_profile.env} 中的 API env key 冲突。
  \item Codex 使用 \cfg{key/url} 时会关闭 inherited config。
  \item 使用 \cfg{key/url} 不能同时开启 inherited auth。
  \item \cfg{provider_profile.home} 目前只支持 Codex runtime home override。
\end{itemize}

\section{maintenance heartbeat}

配置示例：

\begin{lstlisting}
[maintenance.heartbeat]
enabled = true
assessor = "ccb_self"
interval_s = 3600
min_interval_s = 300
unknown_streak_cap = 3
escalation_policy = "report_only"
startup_ensure = true
\end{lstlisting}

默认值：

\begin{itemize}
  \item \cfg{enabled = false}
  \item \cfg{assessor = "ccb_self"}
  \item \cfg{interval_s = 3600}
  \item \cfg{min_interval_s = 300}
  \item \cfg{unknown_streak_cap = 3}
  \item \cfg{escalation_policy = "report_only"}
  \item \cfg{startup_ensure = true}
\end{itemize}

\section{配置验证}

使用：

\begin{lstlisting}
ccb config validate
\end{lstlisting}

如果修改了 \src{.ccb/ccb.config}，先 validate，再启动或 reload。

\section{历史配置补充}

本书以当前 \cfg{version = 2} 的 \src{.ccb/ccb.config} 为主线，不提供 version 1 新项目模板。如果你从旧文档或旧项目看到 version 1 配置，应先迁移到 version 2，再运行 \cmd{ccb config validate}。不要把 version 1 示例复制到新项目中。
